% !TEX root = ../main.tex
% Fonte completa: docs/use-cases/UC-03-2-riattivare-account.md
% Scheda compatta; per i test sono riportate le classi, non tutti i metodi.
\clearpage
\subsection{UC-03.2 - Riattivare un account}
\label{uc:UC-03-2-riattivare-account}
\begingroup
\fontsize{9.2}{10.6}\selectfont
\setlength{\tabcolsep}{5pt}
\setlength{\extrarowheight}{2pt}
\arrayrulecolor{lineeUC}
\setlength{\arrayrulewidth}{0.3pt}
\renewcommand{\arraystretch}{1.12}
\setlength{\parskip}{0pt}
\noindent\begin{tabularx}{\textwidth}{>{\columncolor{etichetteUC}\raggedright\arraybackslash\sffamily\bfseries\color{black!75}}p{2.5cm}>{\raggedright\arraybackslash}X}
\hline
ID e livello & UC-03.2 -- User goal \\[2pt] \hline
Attori & \textbf{Principale:} Administrator; \textbf{secondari:} User \\[2pt] \hline
Descrizione & Un Administrator riattiva un account ordinario precedentemente bannato.\par\textbf{Trigger:} L'Administrator richiede di riattivare uno User bannato. \\[2pt] \hline
Precondizioni & \textbf{1.} Esiste un account Administrator identificabile dal sistema. \par \textbf{2.} Lo User da moderare esiste. \\[2pt] \hline
Postcondizioni & \textbf{1.} Lo User si trova nello stato \texttt{ACTIVE}. \par \textbf{2.} Gli altri dati dello User rimangono invariati. \\[2pt] \hline
Garanzia minima & Se autenticazione, autorizzazione o transizione falliscono, lo stato dello User rimane invariato. \\[2pt] \hline
Flusso principale & \begin{minipage}[t]{\linewidth}\begin{enumerate}[leftmargin=1.4em,nosep,label=\arabic*.]
\item L'Administrator seleziona uno User bannato e richiede di riattivarlo.
\item Il sistema verifica identificativo, password e ruolo dell'Administrator.
\item Il sistema carica lo User indicato.
\item Il sistema verifica che il bersaglio sia uno User ordinario diverso dall'attore.
\item Il sistema verifica che lo User sia \texttt{BANNED}.
\item Il sistema aggiorna atomicamente lo stato a \texttt{ACTIVE}.
\item Il sistema conferma l'operazione.
\end{enumerate}\vspace{3pt}\end{minipage} \\[2pt] \hline
Flussi alternativi & \textbf{2a. Credenziali amministrative non valide} \par (1) Il sistema nega l'operazione. (2) Il caso d'uso termina senza modifiche. \par\smallskip \textbf{3a. User inesistente} \par (1) Il sistema segnala che lo User non esiste. (2) Il caso d'uso termina senza modifiche. \par\smallskip \textbf{4a. Auto-moderazione o bersaglio Administrator} \par (1) Il sistema rileva che l'attore coincide con il bersaglio oppure che il bersaglio è Administrator. (2) Il sistema nega l'operazione. \par\smallskip \textbf{5a. User già attivo} \par (1) Il sistema rileva che lo User è già \texttt{ACTIVE}. (2) Il sistema rifiuta la transizione come conflitto. \\[2pt] \hline
Regole &\par\smallskip\textbullet{} \textbf{BR-05:} le operazioni amministrative richiedono identificativo e password validi di un Administrator.\par\smallskip\textbullet{} \textbf{BR-06:} un Administrator non può moderare se stesso né un altro Administrator. \\[2pt] \hline
Test & \texttt{User\allowbreak{}Service\allowbreak{}Workflow\allowbreak{}Test}. \\[2pt] \hline
\end{tabularx}\par
\endgroup
